% Source: https://artofproblemsolving.com/community/c3h2029806_happy_pi_day_maybe_a_competition
% Stable identity: aops-topic-2029806-q04
The continued fraction of \(\pi\) begins \(\pi=[3;7,15,1,292,\dots]\), with convergents
\[
\dfrac{3}{1},\ \dfrac{22}{7},\ \dfrac{333}{106},\ \dfrac{355}{113},\ \dots
\]
\emph{\(\dfrac{355}{113}\) works:} \(\dfrac{355}{113}=3.14159292\ldots\) while \(\pi=3.14159265\ldots\), so the error is about \(2.67\cdot10^{-7}<3.33\cdot10^{-7}\).

\emph{No smaller denominator works.} By the best-approximation property of convergents, if \(q<113\) (the next convergent denominator after \(106\)), then for every integer \(p\),
\[
|q\pi-p|\ge|106\pi-333|>0.0088 .
\]
Hence for \(q\le112\), \(\left|\pi-\dfrac {p}{q}\right|\ge\dfrac{0.0088}{112}>7.8\cdot10^{-5}\), far larger than \(\dfrac{1}{3}\cdot10^{-6}\).

So the fraction is \(\boxed{\dfrac{355}{113}}\).
